\begin{solution}{16}
\begin{subsolution}
在地面参考系 $S$ 中，从张同学和王同学相遇，到他们之间的距离为 $L$ 时，时间间隔为
\begin{equation}
\Delta t = \frac {L}{2 v}
\label{eq:1.s16}
\end{equation}
在张同学自身静止的参考系 $S'$ 中，张同学开始两次读其随身携带的时钟的两个事件发生在同一地点，因此张同学拍手时其随身携带的时钟的读数 $T$ 为固有时间或原时。利用时间的相对论效应得
\begin{equation}
T = \Delta t \sqrt {1 - \frac {v ^ {2}}{c ^ {2}}} = \Delta t \sqrt {1 - \beta^ {2}} = \frac {L}{2 v} \sqrt {1 - \beta^ {2}}。
\label{eq:2.s16}
\end{equation}
\end{subsolution}

\begin{subsolution}
在参考系 $S'$ 中，张同学第一次拍手这一事件 $1$ 发生的时刻和位置分别是 $t_1' = T$ 和 $x_1' = 0$ 。已设张、王同学分别位于各自参考系中直角坐标系的坐标原点，两参考系中直角坐标系的原点在 $t' = t'' = 0$ 时重合、 $x'$ 和 $x''$ 轴重合且同向。而在王同学自身静止的参考系 $\mathbf{S}''$ 中，张同学的时钟慢了。利用时间的相对论效应，在参考系 $\mathbf{S}''$ 中事件 $1$ 发生的时刻 $t_1''$ (此即王同学随身携带的时钟的读数) 为
\begin{equation}
t _ {1} ^ {\prime \prime} = \gamma t _ {1} ^ {\prime} = \gamma T,
\label{eq:3.s16}
\end{equation}
式中
\begin{equation}
\gamma = \frac {1}{\sqrt {1 - v _ {r} ^ {2} / c ^ {2}}} = \frac {1 + \beta^ {2}}{1 - \beta^ {2}},
\label{eq:4.s16}
\end{equation}
推导中已应用了张同学相对于王同学做匀速直线运动的速率 $v_{r}$ 的表达式。注意：在参考系 $\mathbf{S}''$ 中事件 $1$ 发生的位置并不在王同学自身所在处或原点，而在 $x_{2}^{\prime \prime} \neq 0$ 。

依题意，王同学在时刻
\begin{equation}
t _ {2} ^ {\prime \prime} = t _ {1} ^ {\prime \prime} = \gamma T
\label{eq:5.s16}
\end{equation}
拍手，这一事件在参考系 $\mathbf{S}''$ 中发生的时刻和位置分别是 $t_2'' = t_1''$ 和原点 $x_2'' = 0$ ，是和事件 $1$ 不同的另一事件，记为事件 $2$。而在张同学自身静止的参考系 $S'$ 中，王同学的时钟慢了。利用时间的相对论效应，在参考系 $S'$ 中事件 $2$ 发生的时刻 $t_2'$ 为
\begin{equation}
t _ {2} ^ {\prime} = \gamma t _ {2} ^ {\prime \prime} = \gamma (\gamma T) = \gamma^ {2} T,
\label{eq:6.s16}
\end{equation}
依题意，张同学在时刻
\begin{equation}
t _ {3} ^ {\prime} = t _ {2} ^ {\prime} = \gamma^ {2} T
\label{eq:7.s16}
\end{equation}
第二次拍手(事件 $3$)， $t_{3}^{\prime}$ 即为其随身携带的时钟的读数。

照此继续下去。当张同学第 $n$ 次拍手时，第 $2n - 1$ 个事件在参考系 $\mathbf{S}'$ 中发生的时刻 $t_{2n - 1}'$ (或张同学随身携带的时钟的读数) 为
\begin{equation}
t _ {2 n - 1} ^ {\prime} = \gamma^ {2 (n - 1)} T
\label{eq:8.s16}
\end{equation}
由于在张同学每次观察到的时间间隔 $T$ 内，张、王同学之间在地面参考系中的距离增加 $L$ 。于是，张同学第 $n$ 次拍手时，张、王同学之间在地面参考系中的距离为
\begin{equation}
d _ {n} = \gamma^ {2 (n - 1)} L = \left(\frac {1 + \beta^ {2}}{1 - \beta^ {2}}\right) ^ {2 (n - 1)} L
\label{eq:9.s16}
\end{equation}
\end{subsolution}

\begin{subsolution}
利用时间的相对论效应，从李同学自身静止的参考系看，张同学和王同学依次拍手的时刻为
\begin{equation}
t _ {1} = \left(1 - \frac {v ^ {2}}{c ^ {2}}\right) ^ {- 1 / 2} t _ {1} ^ {\prime} = \left(1 - \beta^ {2}\right) ^ {- 1 / 2} T = \frac {L}{2 v},
\label{eq:10.s16}
\end{equation}
\begin{equation}
t _ {2} = \left(1 - \frac {v ^ {2}}{c ^ {2}}\right) ^ {- 1 / 2} t _ {2} ^ {\prime \prime} = \left(1 - \frac {v ^ {2}}{c ^ {2}}\right) ^ {- 1 / 2} (\gamma T) = \gamma \frac {L}{2 v} = \frac {1 + \beta^ {2}}{1 - \beta^ {2}} \frac {L}{2 v},
\label{eq:11.s16}
\end{equation}
\begin{equation}
t _ {3} = \left(1 - \frac {v ^ {2}}{c ^ {2}}\right) ^ {- 1 / 2} t _ {3} ^ {\prime} = \left(1 - \frac {v ^ {2}}{c ^ {2}}\right) ^ {- 1 / 2} (\gamma^ {2} T) = \gamma^ {2} \frac {L}{2 v} = \left(\frac {1 + \beta^ {2}}{1 - \beta^ {2}}\right) ^ {2} \frac {L}{2 v},
\label{eq:12.s16}
\end{equation}
\begin{equation}
t _ {2 n - 1} = \left(1 - \frac {v ^ {2}}{c ^ {2}}\right) ^ {- 1 / 2} t _ {2 n - 1} ^ {\prime} = \left(1 - \frac {v ^ {2}}{c ^ {2}}\right) ^ {- 1 / 2} (\gamma^ {2 (n - 1)} T) = \gamma^ {2 (n - 1)} \frac {L}{2 v} = \left(\frac {1 + \beta^ {2}}{1 - \beta^ {2}}\right) ^ {2 (n - 1)} \frac {L}{2 v}
\label{eq:13.s16}
\end{equation}
由\eqref{eq:10.s16}至\eqref{eq:13.s16}式可直接看出
\begin{equation}
t _ {1} < t _ {2} < t _ {3} < \dots < t _ {2 n - 1}
\label{eq:14.s16}
\end{equation}
\end{subsolution}
\end{solution}